% The chapter records implementation challenges and the design responses used in this project.
\chapter{Challenges}
\label{ch:challenges}

Despite the progress, we encountered several challenges during the development and evaluation of the forecasting system. They involved production-system representation, supervision and evaluation, temporal dependencies, and multi-task optimization.

\begin{itemize}[leftmargin=1.7em,itemsep=0.55\baselineskip,topsep=0.35\baselineskip]
\item \textbf{Representing heterogeneous factory resources.} Workstations, workers, buffers, AGVs, gantries, and material-supply observations do not share the same semantics or operational role. Treating them as interchangeable features made it difficult to distinguish a processing constraint from a transport or supply constraint. We addressed this by building a 38-node heterogeneous process--logistics graph, using resource-type projections and semantic feature groups, and keeping geographic and semantic relations separate. The NoGraph, NoGroup, NoCross, and NoInfo ablations helped check whether the representation preserved useful cross-resource information.

\item \textbf{Constructing process-informed bottleneck labels.} The main difficulty was converting noisy process signals into labels that reflect a sustained bottleneck. Queueing, blockage, starvation, worker availability, material shortage, and transport delay can overlap, and a rule that works for a workstation may be inappropriate for a gantry or transport robot. We therefore used resource-specific rules, required at least eight minutes of persistence, filled internal gaps of at most one window, and separated Upcoming from Ongoing events. These choices made the labels more consistent with sustained constraints while keeping them computable from history only.

\item \textbf{Learning from rare events under severe class imbalance.} Roughly 98\% of the test pairs are No-event, so a model can appear accurate by favouring the majority class. At the same time, aggressively recovering the few positive windows can produce unstable false alarms. We mitigated this tension with positive-window oversampling, weighted objectives, far-onset weighting, and a hard-false-positive penalty. Event F1, Upcoming/Ongoing recall, and strict onset matching complemented accuracy.

\item \textbf{Capturing delayed cross-resource effects.} Bottleneck evidence does not always appear first at the resource that will later be affected; changes in an upstream station, transport route, or buffer may become visible elsewhere only after several minutes. This made it difficult to distinguish a meaningful delayed dependency from a short-lived local fluctuation. We addressed the problem with temporal-pattern prototypes and delay-aware cross-resource attention, allowing earlier resource trajectories to condition later predictions. Corresponding ablation experiments were used to assess the effectiveness of these design choices.

\item \textbf{Balancing prediction targets with different scales.} The system had to predict station, onset, duration, cause, and remaining order time at the same time, although these targets have different types, scales, and error costs. A single training schedule could allow one objective to dominate the others or could disturb already useful event outputs when the remaining-time objective was refined. We used two-stage training: Stage~1 jointly learns the shared representation and event-related outputs, while Stage~2 freezes the encoder, event heads, and cause head and fine-tunes only the remaining-time head. This separated the optimization pressures while preserving the event and cause predictions.
\end{itemize}
